% !TEX program = xelatex
\documentclass[11pt,a4paper]{ctexart}

\usepackage{amsmath,amssymb,amsfonts}
\usepackage{booktabs}
\usepackage{geometry}
\usepackage{hyperref}
\usepackage{enumitem}
\usepackage{xcolor}
\usepackage{longtable}
\usepackage{array}

\geometry{left=2.5cm,right=2.5cm,top=2.5cm,bottom=2.5cm}
\hypersetup{colorlinks=true,linkcolor=blue,urlcolor=blue}

\newcolumntype{L}[1]{>{\raggedright\arraybackslash}p{#1}}

\title{%
  \textbf{星闪 SLB 物理层}\\[0.4em]
  \Large 6.2.2 基础帧结构和物理资源技术文档\\[0.3em]
  \large CP-OFDM 波形 \textbullet\ 超帧/无线帧 \textbullet\ COT/TTI \textbullet\ 信道载波 \textbullet\ RE 网格
}
\author{依据 T/XS 10001-2025《星闪无线通信系统 接入层\\
同步低时延宽带空口 SLB 技术要求和测试方法》整理\\
（团体标准 20250326 版）}
\date{2026 年 7 月 24 日}

\begin{document}
\maketitle
\tableofcontents
\newpage

%======================================================================
\part{总览}
%======================================================================

\section{文档范围与模块划分}

本文档对应 SLB 120\,kHz 物理层协议第 6.2.2 节``基础帧结构和物理资源''，涵盖 6.2.2.1\textasciitilde 6.2.2.10 全部子节。该节在 6.2.1 内容分类之上，定义 120\,kHz CP-OFDM 系统的\textbf{时频资源坐标系}——波形参数、帧层次、COT/TTI、载波结构、RE 索引及通信域/直通链路资源边界，是 6.2.3\textasciitilde 6.2.5 进行信号/控制/数据映射的前置基础。

\begin{longtable}{L{2.2cm}L{4.8cm}L{6.5cm}}
\caption{文档范围与模块划分} \\
\toprule
\textbf{子节号} & \textbf{子节主题} & \textbf{核心输出/行为} \\
\midrule
6.2.2.1 & CP-OFDM 波形与五种符号类型 & $\Delta f$、CP/GAP 长度、边界符号规则 \\
6.2.2.2 & 超帧/无线帧与链路符号 & 1\,ms 超帧、$N_{\mathrm{sym}}^{\mathrm{frame}}$、G/T/A/D 符号配置 \\
6.2.2.3 & COT 与 TTI & 占信道时长、$2^{N_{\mathrm{TTI}}}$ 无线帧块、链路时分规则 \\
6.2.2.4 & 信道与载波 & 20\,MHz 基础信道、161 子载波、$N_{\mathrm{CH}}$/$L_{\mathrm{CH}}$ \\
6.2.2.5 & 天线端口 & 逻辑端口定义；SAB/G-PCIB/广播同端口 \\
6.2.2.6 & 资源网格 & $161 \times N_{\mathrm{sym}}^{\mathrm{frame}}$ 时频平面 \\
6.2.2.7 & 资源元素 RE & $(k,l)_p$ 索引、$a_{k,l}$ 映射、DC/空闲约束 \\
6.2.2.8 & 通信域 & G 节点工作载波上的可调度资源集合 \\
6.2.2.9 & G/T/附链路 & 通信域内三条逻辑链路的方向与符号类型 \\
6.2.2.10 & 直通链路 & 竞争获取的 D 符号资源；先发/后发节点 \\
\bottomrule
\end{longtable}

\section{与相关章节的关系}

\begin{longtable}{L{2.2cm}L{2.2cm}L{9.1cm}}
\caption{与相关章节的关系} \\
\toprule
\textbf{相关节号} & \textbf{关系} & \textbf{说明} \\
\midrule
6.1 & 上游 & $T_s$、$f_s$、比特序约定；CP-OFDM 波形总体框架 \\
6.2.1 & 上游 & 信号/控制/数据/PIB 分类；6.2.2 提供映射坐标 \\
6.2.3 & 下游 & FTS/STS、DMRS 等在 RE $(k,l)$ 上的映射 \\
6.2.4/6.2.5 & 下游 & SAB/RAB/G-PCIB 块内符号 \#3/\#5 等索引依赖帧结构 \\
6.2.9/6.2.10 & 下游 & OFDM 合成与射频发射使用符号时长与 CP 参数 \\
6.5.3 & 下游 & 符号号 $l$ 与循环前缀类型 $N_{\mathrm{CP}}$ 依赖符号类型与边界 CP 配置 \\
8.3 & 上游 & 基础信道中心频率与多载波频域布局 \\
6.3 & 平行 & 灵活带宽在相同帧结构上的 $\alpha$ 缩放 \\
\bottomrule
\end{longtable}

\section{基本时间与频率量}

\begin{longtable}{llll}
\caption{基本时间与频率量} \\
\toprule
\textbf{基本量} & \textbf{定义/公式} & \textbf{数值} & \textbf{标准条款} \\
\midrule
基准频率 $f_s$ & 物理层采样/定时基准 & 30.72\,MHz & 6.1 \\
基本时间单位 $T_s$ & $T_s = 1/f_s$ & $\approx 32.55$\,ns & 6.1 \\
子载波间隔 $\Delta f$ & $\Delta f = f_s/256$ & 120\,kHz & 6.2.2.1 \\
OFDM 有效部分 & 不含 CP 的数据段 & $256 T_s \approx 8.333$\,$\mu$s & 6.2.2.1 \\
超帧时长 $T_f$ & $T_f = 30720 \times T_s$ & 1\,ms & 6.2.2.2 \\
\bottomrule
\end{longtable}

\newpage
%======================================================================
\part{6.2.2 基础帧结构和物理资源（工程解读）}
%======================================================================

\section{协议章节对应}

本节对应 T/XS 10001-2025 第 6.2 章第 6.2.2 节``基础帧结构和物理资源''，完整涵盖 6.2.2.1\textasciitilde 6.2.2.10 共 10 个子节。

\section{功能定义}

6.2.2 的核心功能为：定义 120\,kHz CP-OFDM 系统的\textbf{时频资源框架}，使发送/接收端能够：
\begin{enumerate}[nosep]
  \item \textbf{确定符号物理参数}：按符号类型查表~\ref{tab:622-symbol-design} 获取 CP/GAP 长度与符号总时长；
  \item \textbf{定位时间资源}：在超帧/无线帧/符号三级时间轴上索引 PIB 与业务资源；
  \item \textbf{索引频域资源}：在节点工作载波上按 $k$、子载波组、工作子载波组编号；
  \item \textbf{映射 RE}：在 $161 \times N_{\mathrm{sym}}^{\mathrm{frame}}$ 资源网格上以 $(k,l)_p$ 写入 $a_{k,l}$；
  \item \textbf{区分域与链路}：通信域（G/T/A 符号）与直通链路（D 符号）及 COT/TTI 占用规则。
\end{enumerate}

\section{模块边界（上下游及职责划分）}

\subsection{上游模块}

\begin{longtable}{L{4cm}L{4.5cm}L{5.5cm}}
\caption{上游模块输入} \\
\toprule
\textbf{上游} & \textbf{输入} & \textbf{说明} \\
\midrule
6.1 物理层一般描述 & $T_s$、$f_s$、比特顺序约定 & 全局 PHY 基准 \\
高层/MAC 配置 & 符号类型、$N_{\mathrm{TTI}}$、$N_{\mathrm{CH}}$、$L_{\mathrm{CH}}$ & 帧长与带宽配置 \\
8.3 信道号与中心频率 & 基础信道中心频率 & 多载波频域布局 \\
\bottomrule
\end{longtable}

\subsection{下游模块}

\begin{longtable}{L{4cm}L{4.5cm}L{5.5cm}}
\caption{下游模块输出} \\
\toprule
\textbf{下游} & \textbf{输出} & \textbf{说明} \\
\midrule
6.2.3 物理层信号 & RE $(k,l)$ 时频位置 & FTS/STS、DMRS、CSI-RS 等映射基准 \\
6.2.4/6.2.5 控制/数据 & 资源网格与符号 \# 索引 & PIB 块内 RE 分配 \\
6.2.9 基带信号生成 & 符号时长与 CP 参数 & OFDM 调制链输入 \\
6.5.3 QPSK 序列 & $n$、$l$、$N_{\mathrm{ID}}$、$N_{\mathrm{CP}}$ & 符号号 $l$ 与 CP 类型对应 \\
\bottomrule
\end{longtable}

\subsection{本模块不负责的内容}

\begin{itemize}[nosep]
  \item 波形调制算法（6.2.9/6.2.10）；
  \item 具体 PIB 比特字段（6.2.4）、信号序列生成（6.2.3）；
  \item 信道编码（6.2.6）、射频指标（第 8 章）；
  \item MAC 调度策略与竞争算法细节。
\end{itemize}

\section{在 SLB 通信流程中的角色}

6.2.2 处于 PHY 层\textbf{资源坐标定义层}：在 6.2.1 内容分类之后、6.2.3\textasciitilde 6.2.5 具体映射之前。

\textbf{接收端资源定位故事线：}
\begin{enumerate}[nosep]
  \item \textbf{帧同步}：锁定超帧/无线帧边界（依赖 6.2.3 FTS/STS，坐标系来自本节）；
  \item \textbf{符号定时}：按符号类型与表~\ref{tab:622-boundary-cp} 区分一般/边界 CP-OFDM 符号；
  \item \textbf{载波选择}：确定节点工作信道/载波及基础载波序号；
  \item \textbf{RE 解调}：在 $(k,l)_p$ 上读取 $a_{k,l}$，区分 G/T/A/D 符号上下文；
  \item \textbf{块解析}：结合 6.2.1 PIB 分类，在已定位 RE 上解析控制/数据。
\end{enumerate}

\textbf{坐标轴关系：}
\begin{itemize}[nosep]
  \item \textbf{时间轴}：超帧 (1\,ms) $\rightarrow$ 8 个无线帧 $\rightarrow$ $N_{\mathrm{sym}}^{\mathrm{frame}}$ 个符号 $\rightarrow$ RE 时域索引 $l$；
  \item \textbf{频域轴}：节点工作信道 ($N_{\mathrm{CH}}$ 个 20\,MHz) $\rightarrow$ 节点工作载波 $\rightarrow$ 161 子载波/基础载波 $\rightarrow$ RE 频域索引 $k$。
\end{itemize}

%----------------------------------------------------------------------
\section{关键参数与强制要求（提炼版）}
%----------------------------------------------------------------------

\subsection{CP-OFDM 波形与符号类型（6.2.2.1）}

物理层信号、控制、数据均基于 $\Delta f = 120$\,kHz CP-OFDM 传输。有效数据段长度恒为 $256 T_s$；不含 CP 的 OFDM 符号满足 $T_{\mathrm{CP}}=0$，$T_{\mathrm{Symb}}^{\mathrm{OFDM}} = 256 T_s$。在无需区分时，OFDM 符号与 CP-OFDM 符号统称``符号''。

\begin{longtable}{L{1.6cm}lllll}
\caption{120\,kHz CP-OFDM 符号设计（提炼，$T_s$ + $\mu$s）\label{tab:622-symbol-design}} \\
\toprule
\textbf{符号类型} & \textbf{一般 CP} & \textbf{一般符号长} &
\textbf{GAP1} & \textbf{边界符号长} & \textbf{边界符号} \\
\midrule
\endfirsthead
\multicolumn{6}{c}{\small（续表）} \\
\toprule
\textbf{符号类型} & \textbf{一般 CP} & \textbf{一般符号长} &
\textbf{GAP1} & \textbf{边界符号长} & \textbf{边界符号} \\
\midrule
\endhead
\bottomrule
\endfoot
类型一 A &
$18T_s$（0.586\,$\mu$s） &
$274T_s$（8.919\,$\mu$s） &
$20T_s$（0.651\,$\mu$s） &
$276T_s$（8.984\,$\mu$s） &
有（\#0,\#7） \\
类型一 B &
$18T_s$（0.586\,$\mu$s） &
$274T_s$（8.919\,$\mu$s） &
$22T_s$（0.716\,$\mu$s） &
$278T_s$（9.049\,$\mu$s） &
有（\#0） \\
类型二 &
$39T_s$（1.270\,$\mu$s） &
$295T_s$（9.603\,$\mu$s） &
$44T_s$（1.432\,$\mu$s） &
$300T_s$（9.766\,$\mu$s） &
有（\#0） \\
类型三 &
$64T_s$（2.083\,$\mu$s） &
$320T_s$（10.417\,$\mu$s） &
--- & --- & 无 \\
类型四 &
$128T_s$（4.167\,$\mu$s） &
$384T_s$（12.500\,$\mu$s） &
--- & --- & 无 \\
\end{longtable}

\textbf{要点}：类型一 A/B 一般 CP 相同，差异在 GAP1/边界符号长度及边界符号位置；类型三/四无单独边界符号与 GAP2 配置。

\textbf{与 6.5.3 衔接}：通信域 CP-OFDM 符号类型 $\rightarrow$ $N_{\mathrm{CP}}$：类型一 A/B $\rightarrow$ 3；类型二 $\rightarrow$ 2；类型三 $\rightarrow$ 1；类型四 $\rightarrow$ 0。

\subsection{超帧/无线帧与符号配置（6.2.2.2）}

\begin{longtable}{lll}
\caption{超帧与无线帧关键参数} \\
\toprule
\textbf{参数} & \textbf{取值/规则} & \textbf{标准条款} \\
\midrule
超帧时长 $T_f$ & $30720 T_s = 1$\,ms & 6.2.2.2 \\
超帧编号 & \#0\textasciitilde\#65535 循环 & 6.2.2.2 \\
每超帧无线帧数 & 8（\#0\textasciitilde\#7） & 6.2.2.2 \\
$N_{\mathrm{sym}}^{\mathrm{frame}}$ & 类型一 A/B: 14；二: 13；三: 12；四: 10 & 表~\ref{tab:622-sym-per-frame} \\
通信域符号 & G/T/A 符号 + 切换间隔 & 6.2.2.2 \\
直通链路符号 & D 符号 + 切换间隔 & 6.2.2.2 \\
\bottomrule
\end{longtable}

\begin{longtable}{lc}
\caption{每无线帧符号数\label{tab:622-sym-per-frame}} \\
\toprule
\textbf{符号类型} & \textbf{$N_{\mathrm{sym}}^{\mathrm{frame}}$} \\
\midrule
\endfirsthead
\multicolumn{2}{c}{\small（续表）} \\
\toprule
\textbf{符号类型} & \textbf{$N_{\mathrm{sym}}^{\mathrm{frame}}$} \\
\midrule
\endhead
\bottomrule
\endfoot
类型一 A / 类型一 B & 14 \\
类型二 & 13 \\
类型三 & 12 \\
类型四 & 10 \\
\end{longtable}

\begin{longtable}{llll}
\caption{无线帧符号配置（边界 CP 位置）\label{tab:622-boundary-cp}} \\
\toprule
\textbf{符号类型} & \textbf{符号数} & \textbf{一般 CP 符号 $l$} & \textbf{边界 CP 符号 $l$} \\
\midrule
\endfirsthead
\multicolumn{4}{c}{\small（续表）} \\
\toprule
\textbf{符号类型} & \textbf{符号数} & \textbf{一般 CP 符号 $l$} & \textbf{边界 CP 符号 $l$} \\
\midrule
\endhead
\bottomrule
\endfoot
类型一 A & 14 & \#1\textasciitilde\#6, \#8\textasciitilde\#13 & \#0, \#7 \\
类型一 B & 14 & \#1\textasciitilde\#13 & \#0 \\
类型二 & 13 & \#1\textasciitilde\#12 & \#0 \\
类型三 & 12 & \#0\textasciitilde\#11 & --- \\
类型四 & 10 & \#0\textasciitilde\#9 & --- \\
\end{longtable}

\subsection{COT 与 TTI（6.2.2.3）}

\begin{longtable}{lll}
\caption{COT 与 TTI 关键参数} \\
\toprule
\textbf{参数/规则} & \textbf{取值/约束} & \textbf{标准条款} \\
\midrule
$N_{\mathrm{TTI}}$ & 0\textasciitilde 6 & 6.2.2.3 \\
TTI 无线帧数 & $2^{N_{\mathrm{TTI}}}$（1/2/4/8/16/32/64） & 6.2.2.3 \\
COT 定义 & 竞争成功后首帧至释放前末帧，整数个无线帧 & 6.2.2.3 \\
边界对齐 & COT/TTI 边界与超帧边界\textbf{无固定关系} & 6.2.2.3 \\
通信域 TTI & 至少含 G 链路且 G 在 TTI 开头 & 6.2.2.3 \\
链路切换 & 不同类型链路间插入 1 符号切换间隔 & 6.2.2.3 \\
附/直通多发节点 & 不同发送节点间亦插入 1 符号切换间隔 & 6.2.2.3 \\
TTI 末符号 & 必须为切换间隔 & 6.2.2.3 \\
\bottomrule
\end{longtable}

\begin{longtable}{L{3cm}L{5.5cm}L{5.5cm}}
\caption{连续 / 非连续传输模式对比} \\
\toprule
\textbf{模式} & \textbf{COT} & \textbf{TTI 与超帧} \\
\midrule
\endfirsthead
\multicolumn{3}{c}{\small（续表）} \\
\toprule
\textbf{模式} & \textbf{COT} & \textbf{TTI 与超帧} \\
\midrule
\endhead
\bottomrule
\endfoot
非连续传输 & 竞争抢占；COT = 整数个无线帧；1 个 COT 含整数个 TTI & COT/TTI 边界与超帧\textbf{无固定}对齐 \\
连续传输 & 无竞争 COT 概念（持续占用） & 超帧含整数个 TTI，或 1 个 TTI 含整数个超帧 \\
\end{longtable}

\subsection{信道、载波与子载波组（6.2.2.4）}

\begin{longtable}{lll}
\caption{信道与载波关键参数} \\
\toprule
\textbf{参数} & \textbf{取值/规则} & \textbf{标准条款} \\
\midrule
基础信道带宽 & 20\,MHz & 6.2.2.4 \\
每基础载波子载波数 & 161；\#80 为 DC，不映射 & 6.2.2.4 \\
有效子载波 & 除 \#80 外均可映射 & 6.2.2.4 \\
基础子载波组 & 每 10 个有效子载波 = 1 组，共 16 组 & 6.2.2.4 \\
分隔子载波 & 频域相邻基础载波间，不映射 & 6.2.2.4 \\
节点工作信道 & 连续 $N_{\mathrm{CH}}$ 个基础信道并集 & 6.2.2.4 \\
工作子载波组编号 & \#0\textasciitilde\#$(\lfloor 16 N_{\mathrm{CH}}/L_{\mathrm{CH}} \rfloor - 1)$ & 6.2.2.4 \\
\bottomrule
\end{longtable}

\begin{longtable}{c*{18}{c}}
\caption{$N_{\mathrm{CH}}$ 与 $L_{\mathrm{CH}}$ 对应关系（共 18 档）\label{tab:622-nch-lch}} \\
\toprule
$N_{\mathrm{CH}}$ &
1 & 2 & 3 & 4 & 5 & 6 & 7 & 8 & 9 & 10 & 11 & 12 & 13 & 14 & 15 & 16 & 25 & 32 \\
\midrule
\endfirsthead
\multicolumn{19}{c}{\small（续表）} \\
\toprule
$N_{\mathrm{CH}}$ &
1 & 2 & 3 & 4 & 5 & 6 & 7 & 8 & 9 & 10 & 11 & 12 & 13 & 14 & 15 & 16 & 25 & 32 \\
\midrule
\endhead
\bottomrule
\endfoot
$L_{\mathrm{CH}}$ &
1 & 2 & 4 & 4 & 8 & 8 & 8 & 8 & 16 & 16 & 16 & 16 & 16 & 16 & 16 & 16 & 32 & 32 \\
\end{longtable}

\textbf{注}：$N_{\mathrm{CH}}=25$ 时工作子载波组编号仍按 $\lfloor 16 N_{\mathrm{CH}}/L_{\mathrm{CH}} \rfloor - 1$ 取上界（向下取整）。

\subsection{天线端口、资源网格与 RE（6.2.2.5\textasciitilde 6.2.2.7）}

\begin{longtable}{lll}
\caption{天线端口、资源网格与 RE} \\
\toprule
\textbf{对象} & \textbf{规则} & \textbf{标准条款} \\
\midrule
天线端口 & 同端口相邻符号信道可互推；逻辑概念非物理天线号 & 6.2.2.5 \\
通信域同端口 & SAB、G-PCIB、广播、G-DMRS-C 同端口 & 6.2.2.5 \\
直通同端口 & 同节点发送的控制/数据/信号同端口 & 6.2.2.5 \\
资源网格 & 每端口 $p$、每基础载波、每无线帧：$161 \times N_{\mathrm{sym}}^{\mathrm{frame}}$ & 6.2.2.6 \\
RE 索引 & $(k,l)_p$，$k{=}0{\ldots}160$，$l{=}0{\ldots}N_{\mathrm{sym}}^{\mathrm{frame}}{-}1$ & 6.2.2.7 \\
DC / 空闲 & $a_{80,l}\equiv 0$；未分配 RE 映射 0 & 6.2.2.7 \\
\bottomrule
\end{longtable}

\subsection{通信域、三链路与直通链路（6.2.2.8\textasciitilde 6.2.2.10）}

\begin{longtable}{lll}
\caption{通信域、三链路与直通链路索引} \\
\toprule
\textbf{对象} & \textbf{规则} & \textbf{标准条款} \\
\midrule
通信域 & G 节点工作载波上 SAB、G-PCIB、广播、可调度资源的集合 & 6.2.2.8 \\
G 链路 & G$\rightarrow$至少一个 T；使用 G 符号 & 6.2.2.9 \\
T 链路 & T$\rightarrow$G；使用 T 符号 & 6.2.2.9 \\
附链路 & T$\rightarrow$T 和/或 G；使用 A 符号 & 6.2.2.9 \\
直通链路 & 竞争获取；先发/后发节点；使用 D 符号 & 6.2.2.10 \\
\bottomrule
\end{longtable}

%----------------------------------------------------------------------
\section{本节核心详解}
%----------------------------------------------------------------------

\subsection{6.2.2.1 CP-OFDM 波形}

\textbf{功能}：为全部物理层内容提供统一时域符号结构（循环前缀 + 有效数据段）。

\textbf{输入/输出}：输入符号类型配置 $\rightarrow$ 输出每符号 CP 长度、GAP1、总时长、是否存在边界符号。

\textbf{上下文}：
\begin{itemize}[nosep]
  \item 通信域/直通链路均使用 120\,kHz SCS；
  \item 类型三/四 CP 更长，适用于更大时延扩展/覆盖场景；
  \item 边界符号用于无线帧首尾（表~\ref{tab:622-boundary-cp}），直接影响 6.5.3 中 $N_{\mathrm{CP}}$ 查表与序列生成；
  \item 类型一 A 的 \#7 边界符号与含 SAB 的典型 TTI 前段布局（见 6.2.1）相关。
\end{itemize}

\subsection{6.2.2.2 超帧与无线帧}

\textbf{功能}：建立周期性时间网格，承载 G/T/A/D 链路符号序列。

\textbf{输入/输出}：输入符号类型 $\rightarrow$ 输出 $N_{\mathrm{sym}}^{\mathrm{frame}}$、符号 \#0\textasciitilde\#$N{-}1$ 及一般/边界 CP 类型。

\textbf{链路符号配置}：
\begin{itemize}[nosep]
  \item 通信域无线帧可灵活配置为 G 符号、T 符号、A 符号与切换间隔；
  \item 直通链路无线帧可灵活配置为 D 符号与切换间隔；
  \item 符号编号在无线帧内连续，切换间隔占用 1 个符号槽位。
\end{itemize}

\textbf{上下文}：6.2.4 中 SAB \#0\textasciitilde\#4、广播 \#5/\#7、G-DMRS-C \#6 等符号索引均引用本节无线帧符号编号。

\subsection{6.2.2.3 COT 与 TTI}

\textbf{功能}：描述一次信道占用内的传输时间块，支撑非连续/连续传输模式调度。

\textbf{输入/输出}：输入 $N_{\mathrm{TTI}}$、传输模式、COT 起止无线帧 $\rightarrow$ 输出 TTI 无线帧数、COT 内 TTI 顺序号、当前帧所属 TTI。

\textbf{强制结构}：
\begin{enumerate}[nosep]
  \item 通信域一个 TTI 至少含 G 链路，且 G 链路在 TTI 开头；
  \item 不同类型链路之间插入长度 1 个符号的切换间隔；
  \item 若含附链路或直通链路，不同发送节点之间亦插入 1 符号切换间隔；
  \item 一个 TTI 的最后一个符号必须为切换间隔。
\end{enumerate}

\textbf{上下文}：竞争接入后 COT 起始无线帧不必对齐超帧；实现调度器须支持任意无线帧边界起始的 COT/TTI。

\subsection{6.2.2.4 信道与载波}

\textbf{功能}：定义频域资源层次，支撑多载波聚合与调度资源指示。

\textbf{频域结构示意}：
\begin{itemize}[nosep]
  \item 单基础载波：\#0\textasciitilde\#79 $|$ \#80(DC) $|$ \#81\textasciitilde\#160；
  \item 除 DC 外每连续 10 个有效子载波 = 1 个基础子载波组，共 16 组（\#0\textasciitilde\#15）；
  \item $N_{\mathrm{CH}}$ 个基础载波聚合为节点工作载波；基础子载波组在工作载波上重新编号为 \#0\textasciitilde\#$(16N_{\mathrm{CH}}-1)$；
  \item 从 \#0 起每连续 $L_{\mathrm{CH}}$ 个基础子载波组组成一个节点工作子载波组；
  \item 频域相邻基础载波之间存在分隔子载波，不承载任何内容（见标准图 51）。
\end{itemize}

\textbf{基础子载波组索引（单载波，跳过 DC）}：有效子载波按频率从低到高编号后，组号 $= \lfloor \mathrm{有效序号}/10 \rfloor$，其中有效序号在 $k{<}80$ 时为 $k$，在 $k{>}80$ 时为 $k-1$。

\textbf{上下文}：6.2.4 DSCI 中``基础载波序号''、``工作子载波组''等字段均基于本节编号；$N_{\mathrm{CH}}$、$L_{\mathrm{CH}}$ 必须按表~\ref{tab:622-nch-lch} 配对。

\subsection{6.2.2.5 天线端口}

\textbf{功能}：定义逻辑天线端口，使同端口上不同符号的信道可相互推断，支撑信道估计复用。

\textbf{强制同端口}：
\begin{itemize}[nosep]
  \item 通信域：SAB、G-PCIB 映射的信号/控制，以及广播信息与 G-DMRS-C，天线端口相同；
  \item 直通链路：同一节点发送的全部物理层信号、控制、数据天线端口相同。
\end{itemize}

\subsection{6.2.2.6\textasciitilde 6.2.2.7 资源网格与 RE}

\textbf{资源网格}：对每一个天线端口、每一个基础载波、每一个无线帧，定义 $161 \times N_{\mathrm{sym}}^{\mathrm{frame}}$ 的时频平面。

\textbf{RE}：网格中每个单元为资源元素，由 $(k,l)_p$ 唯一索引，映射复数值 $a_{k,l}^{(p)}$。不产生混淆时可略去 $p$。

\textbf{强制约束}：
\begin{itemize}[nosep]
  \item 不用于传输的 RE 映射 $a_{k,l}{=}0$；
  \item $a_{80,l}\equiv 0$（中心子载波永不承载内容）。
\end{itemize}

\subsection{6.2.2.8 通信域}

\textbf{功能}：在 G 节点工作载波上划定该 G 节点可管理的物理层资源集合，作为通信域内 G/T/附链路与 PIB 映射的资源边界。

\textbf{定义}：在一个 G 节点工作载波上，下列资源组成的集合称为该 G 节点的\textbf{通信域}；拥有该通信域的 G 节点称为该通信域的 G 节点。

\begin{longtable}{L{4.5cm}L{7cm}L{2.5cm}}
\caption{通信域资源组成\label{tab:622-comm-domain}} \\
\toprule
\textbf{资源类别} & \textbf{说明} & \textbf{典型条款} \\
\midrule
\endfirsthead
\multicolumn{3}{c}{\small（续表）} \\
\toprule
\textbf{资源类别} & \textbf{说明} & \textbf{典型条款} \\
\midrule
\endhead
\bottomrule
\endfoot
同步信息块 SAB & G 节点发送的同步信号 + 同步信息 & 6.2.1、6.2.4 \\
G 链路物理层控制信息块 G-PCIB & G 节点发送的物理层控制（调度/指示等） & 6.2.4 \\
广播信息资源 & G 节点发送广播信息的 RE & 6.2.4 \\
可调度/可配置资源 & G 节点可调度或配置的数据、参考信号等 RE & 6.2.3\textasciitilde 6.2.5 \\
\end{longtable}

\textbf{输入/输出}：输入 G 节点工作载波配置（$N_{\mathrm{CH}}$ 等）$\rightarrow$ 输出通信域资源范围标识（是否含 SAB/G-PCIB/广播/可调度资源）。

\textbf{上下文}：
\begin{itemize}[nosep]
  \item 通信域建立在\textbf{节点工作载波}之上，频域边界由 6.2.2.4 决定；
  \item 时域上由 G/T/A 符号与切换间隔组织（6.2.2.2、6.2.2.3）；
  \item 6.2.1 中通信域 PIB（SAB、RAB、G-PCIB、A-PDIB）及散布资源均落在该集合内；
  \item 实现时须能区分``本通信域 G 节点''与其它节点/直通链路资源，避免跨域误解析。
\end{itemize}

\subsection{6.2.2.9 G 链路、T 链路和附链路}

\textbf{功能}：在同一通信域内定义三条逻辑链路，规定发送/接收方向与对应符号类型，支撑 TTI 内时分复用。

一个通信域由该通信域的 \textbf{G 链路}、\textbf{T 链路}和\textbf{附链路（Auxiliary Link）}组成，均在同一 G 节点工作载波/通信域资源池内时分复用。

\begin{longtable}{L{2.2cm}L{4.2cm}L{2cm}L{5cm}}
\caption{通信域三链路定义\label{tab:622-gta-link}} \\
\toprule
\textbf{链路} & \textbf{方向} & \textbf{符号} & \textbf{典型用途} \\
\midrule
\endfirsthead
\multicolumn{4}{c}{\small（续表）} \\
\toprule
\textbf{链路} & \textbf{方向} & \textbf{符号} & \textbf{典型用途} \\
\midrule
\endhead
\bottomrule
\endfoot
G 链路 & G 节点发送 $\rightarrow$ 至少一个 T 节点接收 & G 符号 & 同步、广播、G-PCIB、G 侧数据、CSI-RS 等 \\
T 链路 & T 节点发送 $\rightarrow$ G 节点接收 & T 符号 & 接入、TCI（HARQ/CQI/SR）、SRS、T 侧数据等 \\
附链路 & T 节点发送 $\rightarrow$ 至少一个 T 节点，或至少一个 T 节点与 G 节点接收 & A 符号 & A-PDIB、附链路数据及绑定 DMRS-D 等 \\
\end{longtable}

\textbf{输入/输出}：输入逻辑链路类型 $\rightarrow$ 输出对应符号类型（G/T/A）及允许的发送/接收角色。

\textbf{时分规则（与 6.2.2.3 一致）}：
\begin{itemize}[nosep]
  \item 同一 TTI/无线帧内通过 G/T/A 符号及切换间隔实现时分；
  \item 不同类型链路之间插入长度 1 个符号的切换间隔；
  \item 若含附链路，不同发送节点的附链路之间亦插入 1 符号切换间隔；
  \item 通信域 TTI 至少含 G 链路，且 G 链路位于 TTI 开头。
\end{itemize}

\textbf{上下文}：三链路\textbf{不属于}直通链路；不得将 D 符号配置到通信域 G/T/A 链路，也不得用 G/T/A 符号承载直通链路业务。

\subsection{6.2.2.10 直通链路}

\textbf{功能}：定义无需预先与 G 节点建立连接、由节点竞争信道获得的直连通信资源，支撑设备间物理层交互。

\textbf{定义}：节点竞争信道获取的、用于与\textbf{未建立连接}的节点通信的资源，称为\textbf{直通链路（Passthrough Link）}。

\begin{longtable}{L{3cm}L{9.5cm}}
\caption{直通链路角色与资源\label{tab:622-passthrough}} \\
\toprule
\textbf{概念} & \textbf{说明} \\
\midrule
\endfirsthead
\multicolumn{2}{c}{\small（续表）} \\
\toprule
\textbf{概念} & \textbf{说明} \\
\midrule
\endhead
\bottomrule
\endfoot
先发节点 & 竞争信道成功的节点；发起直通链路占用 \\
后发节点 & 在直通链路上与先发节点通信的其它节点 \\
符号类型 & D 符号 + 切换间隔（见 6.2.2.2） \\
承载内容 & 物理层信号、控制信息、数据信息（如 SAB、P-PCIB、P-DMRS-C、DMRS-D、数据） \\
资源获取 & 竞争抢占形成 COT；无 G 节点预连接调度前置 \\
\end{longtable}

\textbf{输入/输出}：输入竞争成功后的 COT 与节点角色 $\rightarrow$ 输出 D 符号序列及先发/后发通信上下文。

\textbf{与通信域对比}：

\begin{longtable}{L{2.8cm}L{5.5cm}L{5.5cm}}
\caption{通信域与直通链路对比\label{tab:622-domain-vs-passthrough}} \\
\toprule
\textbf{对比项} & \textbf{通信域} & \textbf{直通链路} \\
\midrule
\endfirsthead
\multicolumn{3}{c}{\small（续表）} \\
\toprule
\textbf{对比项} & \textbf{通信域} & \textbf{直通链路} \\
\midrule
\endhead
\bottomrule
\endfoot
拓扑前提 & 经 G 节点管理；含 G/T/附链路 & 设备直连；无需预建 G 连接 \\
符号类型 & G / T / A 符号 & D 符号 \\
资源获取 & G 节点工作载波上的可调度集合 & 竞争信道获得的 COT 资源 \\
典型 PIB & SAB、RAB、G-PCIB、A-PDIB & SAB、P-PCIB \\
控制复杂度 & 完整控制面（广播、接入、HARQ、CQI 等） & 精简，侧重直通调度与数据 \\
\end{longtable}

\textbf{上下文}：直通链路与通信域资源组织\textbf{不可混用}；实现时按域类型选择符号枚举与 PIB 解析入口（见 6.2.1）。若 TTI 内含直通链路且多发送节点时分，不同发送节点之间插入 1 符号切换间隔（6.2.2.3）。

%----------------------------------------------------------------------
\section{数据类型、长度与维度}
%----------------------------------------------------------------------

\begin{longtable}{llll}
\caption{数据类型、长度与维度} \\
\toprule
\textbf{对象} & \textbf{维度/规模} & \textbf{单位} & \textbf{标准条款} \\
\midrule
子载波索引 $k$ & 0\textasciitilde 160（\#80 恒零） & 子载波 & 6.2.2.7 \\
符号索引 $l$ & 0\textasciitilde $N_{\mathrm{sym}}^{\mathrm{frame}}{-}1$ & 符号 & 6.2.2.2 \\
$N_{\mathrm{sym}}^{\mathrm{frame}}$ & 10/12/13/14 & 符号/无线帧 & 表~\ref{tab:622-sym-per-frame} \\
超帧内无线帧 & 8 & 无线帧/超帧 & 6.2.2.2 \\
超帧编号 & 0\textasciitilde 65535 循环 & --- & 6.2.2.2 \\
资源网格 RE 数 & $161 \times N_{\mathrm{sym}}^{\mathrm{frame}}$ & RE/帧/载波/端口 & 6.2.2.6 \\
基础子载波组/载波 & 16 & 组 & 6.2.2.4 \\
每组有效子载波 & 10 & 子载波 & 6.2.2.4 \\
$N_{\mathrm{CH}}$ & 表~\ref{tab:622-nch-lch} 共 18 档 & 基础信道 & 表~\ref{tab:622-nch-lch} \\
TTI 无线帧数 & $2^{N_{\mathrm{TTI}}}$，$N_{\mathrm{TTI}}\in\{0,\ldots,6\}$ & 无线帧 & 6.2.2.3 \\
$a_{k,l}^{(p)}$ & 复数值 & --- & 6.2.2.7 \\
\bottomrule
\end{longtable}

%----------------------------------------------------------------------
\section{时序关系与触发条件}
%----------------------------------------------------------------------

\begin{enumerate}[nosep]
  \item \textbf{超帧周期}：$T_f = 1$\,ms，编号 \#0\textasciitilde\#65535 循环，与业务调度周期独立；
  \item \textbf{无线帧}：超帧内固定 8 帧，符号类型决定每帧符号数与边界 CP 位置；
  \item \textbf{COT 触发}：非连续模式下竞争信道成功后从首无线帧起算，至所有基础载波释放前的末无线帧（含）；
  \item \textbf{TTI 划分}：COT 或连续模式下每 $2^{N_{\mathrm{TTI}}}$ 无线帧为一个 TTI；通信域 $N_{\mathrm{TTI}}$ 为常数；
  \item \textbf{链路时分}：通信域 TTI 内 G 链路在首段，G/T/A 间以 1 符号切换间隔分隔；
  \item \textbf{TTI 结束}：末符号固定为切换间隔，下一 TTI 或释放信道；
  \item \textbf{直通链路}：竞争获取 COT 后配置 D 符号序列，无 G 节点调度前置。
\end{enumerate}

%----------------------------------------------------------------------
\section{处理步骤}
%----------------------------------------------------------------------

接收侧资源定位流程如下；各步用显式条件推进，不引入状态机。

\textbf{Step 1（符号类型 $\rightarrow$ 每帧符号数）}
\begin{itemize}[nosep]
  \item \textbf{进入}：已从广播/同步信息获得符号类型；
  \item \textbf{动作}：查表~\ref{tab:622-sym-per-frame} 得 $N_{\mathrm{sym}}^{\mathrm{frame}}$；查表~\ref{tab:622-symbol-design} 得 CP/GAP 时长；
  \item \textbf{失败}：符号类型非法则报错，不进入后续定位。
\end{itemize}

\textbf{Step 2（帧同步 $\rightarrow$ 边界/一般符号标记）}
\begin{itemize}[nosep]
  \item \textbf{进入}：已完成超帧/无线帧边界锁定（依赖 6.2.3/6.2.7）；
  \item \textbf{动作}：按表~\ref{tab:622-boundary-cp} 标记每帧内边界 CP 符号与一般 CP 符号；
  \item \textbf{失败}：帧边界未锁定则回退同步流程。
\end{itemize}

\textbf{Step 3（COT/$N_{\mathrm{TTI}}$ $\rightarrow$ 当前 TTI）}
\begin{itemize}[nosep]
  \item \textbf{进入}：已知 COT 起始无线帧（若适用）与 $N_{\mathrm{TTI}}$；
  \item \textbf{动作}：计算 TTI 无线帧数 $2^{N_{\mathrm{TTI}}}$、当前 TTI 顺序号及所含无线帧范围；
  \item \textbf{失败}：当前帧不在任何有效 COT/TTI 内则按空闲处理。
\end{itemize}

\textbf{Step 4（符号配置 $\rightarrow$ 链路类型）}
\begin{itemize}[nosep]
  \item \textbf{进入}：已知本 TTI/无线帧的 G/T/A/D 与切换间隔配置；
  \item \textbf{动作}：将每个符号 $l$ 标注为 G/T/A/D 或切换间隔；校验通信域 TTI 以 G 开头且末符号为切换间隔；
  \item \textbf{失败}：结构不满足强制规则则丢弃该 TTI 解析。
\end{itemize}

\textbf{Step 5（$N_{\mathrm{CH}}$ $\rightarrow$ 目标基础载波）}
\begin{itemize}[nosep]
  \item \textbf{进入}：已知节点工作信道配置与目标基础载波序号；
  \item \textbf{动作}：按表~\ref{tab:622-nch-lch} 校验 $N_{\mathrm{CH}}$/$L_{\mathrm{CH}}$，计算工作子载波组编号范围；
  \item \textbf{失败}：表~\ref{tab:622-nch-lch} 配对失败或载波序号越界则报错。
\end{itemize}

\textbf{Step 6（$(k,l)_p$ $\rightarrow$ 解调入口）}
\begin{itemize}[nosep]
  \item \textbf{进入}：已选定端口 $p$、基础载波与无线帧；
  \item \textbf{动作}：索引 RE，校验 $k{\neq}80$ 且目标 RE 非空闲后，进入 6.2.3/6.2.4 解调；
  \item \textbf{失败}：DC/空闲/越界 RE 不得当作有效载荷。
\end{itemize}

%----------------------------------------------------------------------
\section{约束与实现要点}
%----------------------------------------------------------------------

\begin{enumerate}[nosep]
  \item \textbf{【DC 子载波】}：$a_{80,l}\equiv 0$，任何信号/控制/数据均不得映射至 \#80。
  \item \textbf{【分隔子载波】}：多载波频域相邻时，分隔子载波不映射任何内容。
  \item \textbf{【边界 CP 位置】}：必须严格按表~\ref{tab:622-boundary-cp} 配置，否则符号定时与 6.5.3 的 $l$/$N_{\mathrm{CP}}$ 不一致。
  \item \textbf{【TTI 结构】}：通信域 TTI 末符号必须为切换间隔；G 链路必须在 TTI 开头。
  \item \textbf{【边界不对齐】}：COT/TTI 起始可与超帧任意对齐，实现须支持任意无线帧边界。
  \item \textbf{【$N_{\mathrm{CH}}$/$L_{\mathrm{CH}}$ 配对】}：$N_{\mathrm{CH}}$ 与 $L_{\mathrm{CH}}$ 必须按表~\ref{tab:622-nch-lch} 取值，否则工作子载波组编号错误。
  \item \textbf{【空闲 RE】}：未分配 RE 映射 $a_{k,l}{=}0$，接收端不得当作有效载荷。
  \item \textbf{【天线端口】}：通信域 SAB/G-PCIB/广播/G-DMRS-C 须使用同一天线端口做信道估计。
  \item \textbf{【域不可混用】}：通信域 G/T/A 符号配置不适用于直通链路 D 符号场景。
  \item \textbf{【OFDM 简称】}：无 CP 的 OFDM 符号 $T_{\mathrm{CP}}{=}0$；与 CP-OFDM 统称``符号''时须区分上下文。
  \item \textbf{【无状态机】}：帧/TTI/链路/RE 定位用显式条件与查表推进（符号类型、$N_{\mathrm{TTI}}$、域类型），不维护 FSM 跳转表，避免异常路径锁死。
\end{enumerate}

%----------------------------------------------------------------------
\section{与标准其他章节的衔接}
%----------------------------------------------------------------------

\begin{itemize}[nosep]
  \item \textbf{6.2.1}：PIB 与内容分类；6.2.2 提供映射坐标；
  \item \textbf{6.2.3}：FTS/STS/DMRS 等在 RE $(k,l)$ 上的映射规则；
  \item \textbf{6.2.4/6.2.5}：控制/数据在 G/T/A/D 符号上的块内布局与符号 \# 索引；
  \item \textbf{6.5.3}：符号号 $l$ 与 $N_{\mathrm{CP}}$ 依赖表~\ref{tab:622-symbol-design}/表~\ref{tab:622-boundary-cp} 符号类型；
  \item \textbf{6.2.9/6.2.10}：基带 OFDM 合成与射频发射使用本节符号时长参数；
  \item \textbf{8.3}：基础信道中心频率与多载波频域布局；
  \item \textbf{6.3}：灵活带宽在相同帧结构上的带宽扩展。
\end{itemize}

\vfill
\begin{center}
\small\textcolor{gray}{字段级定义与完整参数以 T/XS 10001-2025 第 6 章及 6.2.2 各子节为准；本文档提供 120\,kHz 帧结构与 RE 坐标系工程解读，供 SLB 物理层实现与联调参考。}
\end{center}

\end{document}
